\bookchapter{2}{定态薛定谔方程}

\problem[$\ast$]{2.1}
证明下列三个定理：
\begin{enumerate}
\item 对归一化的解，其分离变量常数 $E$ 必为实数。提示：把 $E$（式（2.7）中）写成 $E_0+\mathrm{i}\Gamma$ 的形式（$E_0$ 和 $\Gamma$ 都是实数），然后证明对任何时间 $t$，如果式（1.20）都成立，则 $\Gamma$ 必定为零。
\item 定态波函数 $\psi(x)$ 总可以取作实数（不像 $\Psi(x,t)$ 一定是复数）：这并不是说任何定态薛定谔方程的解一定是实数，而是说，如果你没有得到实数解，总可以通过这些解（具有相同能量）的线性组合得到一个实数解。所以，你总可以说薛定谔方程的解可以取作实数。提示：对于一个给定的能量 $E$，如果 $\psi(x)$ 满足式（2.5），那么 $\psi(x)$ 的共轭复数也满足；这样它们的线性组合 $\psi+\psi^*$ 和 $\mathrm{i}(\psi-\psi^*)$ 为实数解，它们同样也满足式（2.5）。
\item 如果 $V(x)$ 是偶函数（even function）（即 $V(-x)=V(x)$），那么 $\psi(x)$ 总可以取偶函数或者奇函数。提示：对于给定的能量 $E$，如果 $\psi(x)$ 满足式（2.5），那么 $\psi(-x)$ 也一定满足，所以它们的奇函数和偶函数的组合 $\psi(x)\pm\psi(-x)$ 也满足。
\end{enumerate}
\vfill

\problem[$\ast$]{2.2}
证明对于定态薛定谔方程的每一个归一化解，$E$ 必须大于 $V(x)$ 的最小值。在经典力学中，这句话对应的是什么？提示：把式（2.5）重写为
\[
\frac{\dd^2\psi}{\dd x^2}=\frac{2m}{\hbar^2}[V(x)-E]\psi;
\]
如果 $E<V_{\min}$，那么 $\psi$ 和它的二阶导数有相同的符号，论证该波函数是不可归一化的。
\vfill

\problem{2.3}
对一维无限深方势阱来说，证明在 $E=0$ 或 $E<0$ 的情况下（定态）薛定谔方程不存在可接受的（物理）解。（这是在习题 2.2 中所讨论的一般定理的一个特殊情况，但这次你要直接求解薛定谔方程，并证明其无法满足边界条件。）
\vfill

\problem[$\ast$]{2.4}
对一维无限深势阱中的第 $n$ 个定态，计算 $\avg{x}$、$\avg{x^2}$、$\avg{p}$、$\avg{p^2}$、$\sigma_x$ 和 $\sigma_p$ 的值。验证不确定原理是成立的。指出哪个状态最接近不确定原理的极限？
\vfill

\problem[$\ast$]{2.5}
一维无限深方势阱中粒子初始波函数由前两个定态波函数叠加而成：
\[
\Psi(x,0)=A[\psi_1(x)+\psi_2(x)].
\]
\begin{enumerate}
\item 归一化 $\Psi(x,0)$（即求出 $A$，如果利用 $\psi_1$ 和 $\psi_2$ 的正交归一性会使计算变得很简单。请记住，你可以认为在 $t=0$ 时波函数 $\Psi$ 是归一化的，那么在以后时间也是归一化的——如对此有疑问，在做完（b）后验证一下）。
\item 求 $\Psi(x,t)$ 和 $|\Psi(x,t)|^2$。像例题 2.1 一样，把后者用时间正弦函数表示。为简化结果起见，令 $\omega\equiv\pi^2\hbar/(2ma^2)$。
\item 计算 $\avg{x}$ 的值。注意结果是随时间振荡的。振荡的角频率是多少？振幅是多少？（如果你得到的振幅大于 $a/2$，计算一定有错。）
\item 计算 $\avg{p}$ 的值。（正如彼得·洛所讲：“当然要做，约翰尼！”）
\item 如果你对粒子的能量进行测量，可能得到的值是什么？所获得每个值的几率是多少？求出 $H$ 的期望值，并与 $E_1$ 和 $E_2$ 做比较。
\end{enumerate}
\vfill

\problem{2.6}
虽然波函数的整体相位常数都没有任何物理意义（在计算一个可观测量时可以将此抵消掉），但是在式（2.17）中系数的相对相位却起作用。例如，假定改变习题 2.5 中的 $\psi_1$ 和 $\psi_2$ 的相对相位为
\[
\Psi(x,0)=A[\psi_1(x)+\mathrm{e}^{\mathrm{i}\phi}\psi_2(x)],
\]
其中 $\phi$ 是一些常数。求解 $\Psi(x,t)$、$|\Psi(x,t)|^2$ 和 $\avg{x}$，并与前面（上题）的结果做比较。研究当 $\phi=\pi/2$ 和 $\phi=\pi$ 时的特殊情况。（对该问题的一个图示，可在脚注 9 所给出的网页中找到。）
\vfill

\problem[$\ast$]{2.7}
一维无限深势阱中的粒子初始波函数是
\[
\Psi(x,0)=\begin{cases}
Ax,&0\leq x\leq a/2,\\
A(a-x),&a/2\leq x\leq a.
\end{cases}
\]
\begin{enumerate}
\item 画出 $\Psi(x,0)$ 的草图，然后求出常数 $A$。
\item 求出 $\Psi(x,t)$。
\item 对粒子能量进行测量，得到结果为 $E_1$ 的几率是多少？
\item 利用式（2.21），求出能量的期望值。
\end{enumerate}
\vfill

\problem{2.8}
处在一维无限深（宽度为 $a$）方势阱中质量为 $m$ 的粒子，其初态为
\[
\Psi(x,0)=\begin{cases}
A,&0\leq x\leq a/2,\\
0,&a/2\leq x\leq a.
\end{cases}
\]
对于某一常数 $A$，势阱左半边的区域中（$t=0$ 时）每一处发现粒子的可能性相同。（在以后的时刻 $t$）对能量进行测量得到数值为 $\pi^2\hbar^2/(2ma^2)$ 的几率是多少？
\vfill

\problem{2.9}
对例题 2.2 中的波函数，用公式：
\[
\avg{H}=\int\Psi(x,0)^*\,\hat H\Psi(x,0)\,\dd x
\]
求在 $t=0$ 时 $H$ 的期望值。同在例题 2.3 中用式（2.21）求出的结果做比较。注意：因为 $\avg{H}$ 是不依赖时间的，所以用 $t=0$ 求解也不失其一般性。
\vfill
